 %===============================================================================================================%
%                   Estilo y paquetes que se van a utilizar: No es necesario editar                             %
%===============================================================================================================%
\documentclass[12pt,letterpaper]{article}     % Tipo de documento y otras especificaciones
\usepackage[utf8]{inputenc}                   % Para escribir tildes y eñes
\usepackage[spanish]{babel}                   % Para que los títulos de figuras, tablas y otros estén en español
\addto\captionsspanish{\renewcommand{\tablename}{Tabla}}					% Cambiar nombre a tablas
\addto\captionsspanish{\renewcommand{\listtablename}{Índice de tablas}}		% Cambiar nombre a lista de tablas
\usepackage{geometry}                         
\geometry{left=18mm,right=18mm,top=21mm,bottom=21mm}	 % Tamaño del área de escritura de la página
\usepackage{ucs}
\usepackage{amsmath}      % Los paquetes ams son desarrollados por la American Mathematical Society
\usepackage{amsfonts}     % y mejoran la escritura de fórmulas y símbolos matemáticos.
\usepackage{tikz}

\usetikzlibrary{babel}               % Se añade para que se acepten algunos caractéres en el circuitikz 
\usepackage{circuitikz}
\usepackage{graphicx}     % Para insertar gráficas
\usepackage[lofdepth,lotdepth]{subfig}	% Para colocar varias figuras
\usepackage{unitsdef}	  % Para la presentación correcta de unidades
\usepackage{pdfpages}     % Incluir paginas de pdf externo, para los anexos
\usepackage{appendix}     % Para los anexos
\renewcommand{\unitvaluesep}{\hspace*{4pt}}	% Redimensionamiento del espacio entre magnitud y unidad
\usepackage[colorlinks=true,urlcolor=blue,linkcolor=black,citecolor=red]{hyperref}     % Para insertar 				hipervínculos y marcadores
\usepackage{float}		% Para ubicar las tablas y figuras justo después del texto
\usepackage{booktabs}	% Para hacer tablas más estilizadas
\batchmode
\usepackage{apalike}
\bibliographystyle{apalike}  % Bibliografía en formato IEEE
\pagestyle{plain} 
\pagenumbering{arabic}
\usepackage{lastpage}
\usepackage{skull}
\usepackage{marvosym}
\usepackage{fancyhdr}	% Para manejar los encabezados y pies de página
\pagestyle{fancy}		% Contenido de los encabezados y pies de pagina
\usepackage{multicol}   % Para varias columnas
\usepackage{enumerate}
\usepackage{url}
\usepackage{hyperref}
\usepackage{listings}
\usepackage{xcolor}

\lstset{
  language=C++,
  basicstyle=\ttfamily\footnotesize,
  keywordstyle=\color{blue}\bfseries,
  commentstyle=\color{gray}\ttfamily,
  stringstyle=\color{red!70!black},
  numbers=left,
  numberstyle=\tiny\color{gray},
  stepnumber=1,
  numbersep=5pt,
  backgroundcolor=\color{gray!10},
  frame=single,
  rulecolor=\color{black},
  captionpos=b,
  breaklines=true,
  breakatwhitespace=false,
  tabsize=2,
  showstringspaces=false
}

%===============================================================================================================%
%                         Definición del ambiente: Resumen --> No es necesario modificar                        %                                        
%===============================================================================================================%
\newcounter{resumen}
\setcounter{resumen}{0}
\def\theejemplo{\thechapter.\arabic{resumen}}

\newenvironment{resumen}
{	
	\begin{center}
	\begin{minipage}[t]{500 pt}
	\vspace{1mm}
	\emph{\textbf{}}
	\\[-2mm]
	\line(1,0){500}
	\\[-4.25 mm]
	\line(1,0){500}
	\\
}
{
	\normalsize
	\\[2mm]
	\footnotesize\textbf{Palabras clave: \footnotesize\@palabras}
	\\[-2mm]
	\line(1,0){500}
	\\[0.5cm]
	\end{minipage}
	\end{center}
}

%===============================================================================================================%
%                         Definición del ambiente: Palabras clave --> No es necesario modificar                 %                                                               
%===============================================================================================================%
\def\palabras#1{\gdef\@palabras{#1}}

%===============================================================================================================%
%             Definición de encabezados --> Modificar de acuerdo al número de laboratorio (o materia)           %                                                                          
%===============================================================================================================%
\lhead{IoT} % Codigo de la materia de circuitos 1
\chead{}
\rhead{Manual de Usuario}	% Aquí va el número de Laboratorio, al igual que en el titulo
\lfoot{Programa de Ingeniería Electrónica y Sistemas}
\cfoot{\thepage\ de \pageref{LastPage}}
\rfoot{Universidad Católica de Oriente}

%===============================================================================================================%
% Portada --> Modificar de acuerdo al número de laboratorio, materia, grupo de trabajo, docente, estudiantes y  %    %                                       respectivas identificaciones                                            %                                                                          
%===============================================================================================================%
\author{Santiago Barbarán, 1001159841\\Miguel Ángel Botero,1023523241 \\ Sebastián Moreno Vargas, 1001724352\\Profesor: Juan Camilo Acevedo \vspace*{2.8in}} % Pendiente de modificar el grupo y/o agregar compañero
\title{
\begin{minipage}{\textwidth}
\centering
	\includegraphics[width=0.3\textwidth]{uco.png}
	\bigskip
\end{minipage}\\
Universidad Católica de Oriente \\{\small Facultad de Ingeniería\\Programa de Ingeniería Electrónica y Sistemas\\ Informe IoT\\Semestre 2026-II\\\vspace*{0.55in} Informe:}\\  Manual de Usuario \vspace*{1.35in}}

\begin{document}	% Inicio del documento

\pdfbookmark[1]{Portada}{portada} 	% Marcador para el título, no editar

\maketitle							% Título, no editar
\thispagestyle{empty}

\tableofcontents         % Tabla de contenidos, no editar, se crea automáticamente
%\newpage
%\listoffigures           % Lista de figuras, no editar, se crea automáticamente
%\listoftables            % Lista de tablas, no editar, se crea automáticamente

\newpage

%===============================================================================================================%
%                            Resumen --> Introducir resumen del laboratorio                                     %                                                                          
\section{Introducción y descripción general}

Bienvenido al Manual de Usuario del sistema de monitoreo y control autónomo de calidad de agua para piscicultura. Este sistema ha sido diseñado para supervisar de forma continua y automática los parámetros del agua de los estanques (nivel y calidad mediante sólidos disueltos totales / conductividad), buscando tener condiciones idóneas para la supervivencia y salud de la biomasa de peces.

El sistema funciona de forma autónoma realizando acciones preventivas (activación de bomba de recirculación y electroválvula de recambio hídrico) y permitiendo a los operadores monitorear el estado del estanque desde una interfaz web accesible.\\

\section{Acceso a la plataforma de monitoreo}

Para consultar el estado del estanque y visualizar métricas históricas, siga estos pasos:\\

\begin{itemize}
    \item \textbf{Conexión a la red local:} Asegúrese de que su dispositivo esté conectado a la red local del criadero.

    \item \textbf{Ingreso al Dashboard:} Abra un navegador y en la barra de direcciones, ingrese la dirección IP asignada al servidor del sistema seguido del puerto 3000.
    \item \textbf{Credenciales:} Ingrese su usuario y contraseña asignados por el administrador para acceder al panel principal de Grafana.
\end{itemize}

\section{Interpretación de la Dashboard}
\begin{itemize}
\begin{figure}[H]
    \centering
    \includegraphics[width=0.5\linewidth]{SDT.png}
    \caption{Historial SDT}
\end{figure}

    \item \textbf{Gráfica de Historial (SDT):} En esta gráfica podrá ver la tendencia y variaciones continuas de la concentración de sólidos disueltos a lo largo del tiempo, lo que le permitirá identificar cambios bruscos o picos en la calidad del agua.


\begin{figure}[H]
    \centering
    \includegraphics[width=0.5\linewidth]{Profundidad.png}
    \caption{Historial Profundidad}
    \label{fig:placeholder}
\end{figure}
\item \textbf{Gráfica de Historial (Profundidad):}En este apartado podrá consultar el comportamiento del nivel del estanque a lo largo de la jornada. Le servirá para detectar caídas de nivel, descensos drásticos o confirmar cuando los ciclos automáticos de vaciado y llenado estén operando.

\end{itemize}
\section{Operación del sistema y lógica autónoma}
El sistema EcoSmart UCO no requiere intervención manual constante para regular el agua. Su lógica interna actúa de la siguiente manera:

\begin{enumerate}
    \item \textbf{Recirculación Automática:} Si la calidad del agua o el oxígeno disuelto se salen de los intervalos establecidos, el sistema enciende automáticamente la bomba de recirculación de 110V para oxigenar el medio.
    \item \textbf{Ciclos de Vaciado e Ingreso}: Si los valores de conductividad/SDT superan el umbral máximo tolerado por los peces, se acciona las electroválvulas para drenar un porcentaje del agua acumulada e ingresar agua limpia.
\end{enumerate}

\section{Resolución de problemas frecuentes}
\begin{enumerate}
 


            \item \textbf{Problema:} El Dashboard no carga en el navegador.
            \begin{itemize}
                \item \textbf{Causa Probable:} Apagado de la Raspberry Pi 5.
                \item \textbf{Solución:} Verifique que el LED de encendido de la Raspberry Pi 5 esté activo y fijo.
                
            \end{itemize}
            
            \vspace{0.5cm}
            
            \item \textbf{Problema:} El nivel de agua muestra un valor estático o erróneo.
            \begin{itemize}
                \item \textbf{Causa Probable:} Obstrucción o suciedad en la sonda ultrasónica).
                \item \textbf{Solución:} Verifique que la sonda montada sobre el estanque no tenga objetos, algas ni suciedad en la superficie cónica de medición.
            \end{itemize}
            
            \vspace{0.5cm}
            
            \item \textbf{Problema:} Las lecturas de conductividad oscilan bruscamente.
            \begin{itemize}
                \item \textbf{Causa Probable:} Electrodo sucio o sedimentación en las puntas de prueba del sensor.
                \item \textbf{Solución:} Limpie suavemente las puntas del sensor utilizando agua destilada o agua limpia, evitando frotar con materiales abrasivos.
            \end{itemize}
            
            \vspace{0.5cm}
            

\end{enumerate}

\section{Precauciones de seguridad operacional}
\begin{itemize}
    \item Seguridad Eléctrica: No manipule los gabinetes metálicos ni la caja de los contactores/relés de 110V con las manos mojadas o descalzo.

    \item Limpieza de Sensores: Realice inspecciones visuales semanales a los sensores en el estanque para retirar hojas, ramas o acumulación de algas.


\end{itemize}

\end{document} 

